\documentclass[11pt]{article}
\usepackage{fontspec,amsmath,amssymb,amsthm,hyperref}
\usepackage[margin=1in]{geometry}
\setlength{\emergencystretch}{3em}
\title{Identify the implemented exponential-step error with its actual integral defect}
\author{LangevinLean mathematical exposition}
\date{}
\begin{document}\maketitle
\noindent Source module: \texttt{KineticLangevin/MildError.lean}. Compile this proof document with XeLaTeX or LuaLaTeX. The accompanying HTML pairs the mathematical argument with the exact Lean statements and records the scope of the formalization.


The local errors are concrete integrals of changes in the force. This module proves that these integrals really are the errors of the implemented exponential step when compared with a mild exact velocity and its integrated position. The variation-of-constants equation is an explicit input here; constructing its Brownian convolution and deriving it from the original integral equation remain distinct steps.


\section{Separate the frozen force from its variation}


Write \(D_v(t)=\int_0^t e^{-\gamma(t-s)}[g(x_s)-g(x_0)]ds\). Linearity and \(\int_0^t e^{-\gamma(t-s)}ds=\psi_1(t)\) give

\[D_v(t)=\int_0^t e^{-\gamma(t-s)}g(x_s)ds-\psi_1(t)g(x_0).\]

Lean proves interval integrability using continuity of the path, drift and exponential kernel; it then distributes scalar multiplication and the integral over subtraction. For a mild exact velocity \(V_t=\psi_0(t)V_0-\int_0^t e^{-\gamma(t-s)}g(x_s)ds+B_t\), subtracting it from the frozen velocity plus the same convolution \(B_t\) gives exactly \(D_v(t)\). The proof is algebra after the integral identity.



\noindent Lean declarations: \texttt{velocityDefect\_eq\_integral\_sub}, \texttt{frozenVelocity\_error}.

\section{Integrating the frozen velocity gives the position formula}


The frozen position is \(x+\psi_1(t)v-\psi_2(t)g\). Its derivative is the frozen velocity \(\psi_0(t)v-\psi_1(t)g\), and its value at zero is \(x\). The fundamental theorem of calculus gives \(\int_0^t V_{\rm frozen}(s)ds=X_{\rm frozen}(t)-x\). Lean invokes the already proved derivative theorem and proves integrability from continuity of the coefficients. The identity holds also at zero friction because the coefficients were defined by integrals.



\noindent Lean declarations: \texttt{integral\_frozenVelocity}.

\section{Both coordinates use the same noise convolution}


The position innovation is \(\int_0^h B_sds\), the time integral of the same velocity convolution. If \(x_h-x_0=\int_0^hV_sds\), integrate the velocity error identity to obtain the position defect \(D_x(h)=\int_0^hD_v(s)ds\). Thus

\[\operatorname{exponentialStep}(V_0,x_0;B_h,\int_0^hB_sds)-(V_h,x_h)=(D_v(h),D_x(h)).\]

Lean proves the velocity identity at every point in the interval, changes the integrand using that pointwise equality, distributes the integral over sums and differences, and uses the frozen-velocity integral and the kinematic identity. Extensionality then joins the two coordinate equalities into the phase-space equality. No stochastic error estimate is hidden inside the definitions.



\noindent Lean declarations: \texttt{exponentialStep\_error\_identity}.

\section{Bound the implemented step in the transformed norm}


Apply the exact step-error identity to replace \texttt{exponentialStep − exact endpoint} by the explicit pair \((D_v,D_x)\). The local-error module has already integrated the kinematic velocity bound to prove \(\mathcal E(D_v,D_x)\le(3MC h^2/4)^2\) under \(\gamma h\le1/4\). The norm-square identity and nonnegativity allow \texttt{transformNorm\_le\_of\_energy\_le} to take square roots, giving \(\|\text{step}-\text{exact}\|_T\le3MC h^2/4\). Every input is visible: the actual mild velocity equation, continuous common convolution noise, the kinematic identity, and the uniform velocity norm bound. No stochastic moment or stationary law is asserted here.



\noindent Lean declarations: \texttt{exponentialStep\_localNorm\_bound}.

\end{document}
